% ===================================================================
% SECTION 2: STUDY AREA & DATA SETS
% File: sections/dataset002.tex
% ===================================================================

\subsection{Datasets}
\label{subsec:datasets}

This study integrated four primary data sources to characterize compaction within individual depth sections, comprising (1) multilayer compaction observations, (2) groundwater level records, (3) borehole lithological profiles, and (4) vertical surface displacement derived from Global Navigation Satellite System (GNSS) stations and Small Baseline Subset Interferometric Synthetic Aperture Radar (SBAS-InSAR) analysis. The following subsections describe each dataset, its role in the compaction analysis, and the preprocessing applied before model construction.

\subsubsection{Multilayer aquifer-system compaction}

Multilayer compaction monitoring wells (MLCWs) recorded deformation at multiple depths. These specialized borehole extensometers contained 21 to 26 magnetic rings installed along the well profile to depths of up to 300~m, allowing compaction within individual aquifer and aquitard units to be isolated with a measurement precision of 1~mm \citep{hung_measuring_2021}. Monthly raw observations from five primary MLCW stations (Guangfu, Huwei, Tuku, Honglun, and Xiutan) were collected and partitioned into six uniform 50~m depth sections (designated as S1 through S6, spanning 0--300~m depth) to account for differences in anchor ring depth among boreholes, with station coordinates and monitoring specifications summarized in \Cref{tab:mlcw_info}. Sectional time series were then fitted with a deformation model comprising linear, periodic, and step functions (\Cref{subsec:parametric_deformation}), which estimated linear velocities and annual and semiannual components. Differences between consecutive months were subsequently calculated to obtain compaction increments for each section.

\subsubsection{Groundwater level observations}

Groundwater level (GWL) observations recorded variations in hydraulic head associated with aquifer-system compaction and land subsidence in alluvial fan aquifer systems \citep{galloway_land_1999, gambolati_2015, herrera-garcia_mapping_2021, bagheri-gavkoshLandSubsidenceGlobal2021}. The regional GWL monitoring network operated by the Water Resources Agency (WRA) of Taiwan comprised nested observation wells screened at discrete depth intervals corresponding to Aquifers 1 through 4 \citep{survey_project_1999, liu_characterization_2004}. Daily piezometric head measurements, expressed as elevations relative to Mean Sea Level (m MSL), spanned from January 2000 to December 2025 across Yunlin County, capturing seasonal agricultural pumping drawdowns, monsoon recharge during the wet season, and head declines over multiple decades \citep{chang2022_wetanddry, lu2020_crfp}. Station locations, screened intervals, and monitoring periods are summarized in \Cref{tab:gwl_info}. Daily piezometric head observations were averaged to monthly values, and differences between consecutive months were calculated to obtain monthly hydraulic head changes (\Cref{subsec:hydrogeological_drivers}).

\subsubsection{Borehole lithological profiles}

Borehole lithological profiles described variations in sediment composition among the six 50~m compaction sections. Detailed logs at each MLCW recorded the vertical distribution of gravel, sand, silt, clay, and mixed fine-grained deposits to depths of 300~m. The logged materials were classified as gravel, coarse sand, fine sand, or a combined clay, silt, and mud group within each depth section. These borehole-derived compositions provided the sediment predictors used for model calibration and cross-station validation after the log ratio balance transformation described in \Cref{subsec:ilr_sbp_transformation}.

Sediment architecture between monitoring stations was represented separately by the 3D hydrogeological model developed by the Geological Survey and Mining Management Agency \citep{gsmma_3d}. % TODO(CITATION): Find a peer-reviewed source supporting the 1 m vertical resolution and 500 m horizontal grid spacing of the 3D hydrogeological model; add its BibTeX entry to writing_manu2.bib, then insert \citep{citation_key} here.
The regional model provided continuous lithological profiles at a 1~m vertical resolution and a 500~m horizontal grid spacing. This regional product was retained to describe sediment variability between monitoring stations and did not replace the borehole-derived sediment compositions used in the reported station-based analyses.

\subsubsection{Vertical surface displacement}

Vertical surface displacement measurements recorded the integrated ground displacement resulting from compaction across all underlying depth intervals. Daily 3D position time series from continuous GNSS (cGNSS) stations located at or near the MLCW sites \citep{IESAS_TGM_2026} provided point observations of vertical displacement from 2010 to 2024, with station coordinates and monitoring periods summarized in \Cref{tab:gnss_info}. Surface locations outside GNSS coverage were monitored using Small Baseline Subset (SBAS) InSAR analysis of Sentinel-1 Synthetic Aperture Radar (SAR) imagery, comprising 266 ascending path 69 images and 264 descending path 105 images with acquisition parameters summarized in \Cref{tab:sentinel1_info} \citep{torres_gmes_2012, yague-martinez_interferometric_2016}. SBAS interferogram stacks were generated with the HyP3 pipeline \citep{hogenson_hybrid_2025}, and multitemporal phase inversion was conducted using MintPy \citep{yunjun_small_2019}. These processing steps produced initial displacement measurements along the radar line of sight. Vertical surface displacement was then resolved by vector decomposition in two dimensions \citep{fuhrmann_resolving_2019, hanssen_radar_2001}. Vertical surface displacement time series from cGNSS and SBAS-InSAR were subsequently fitted with a deformation model comprising linear, periodic, and step functions (\Cref{subsec:parametric_deformation}), resampled to a monthly interval, and differenced between consecutive months to obtain surface displacement increments.

Integration of these four preprocessed data sources produced a monthly dataset aligned in space and time. The dataset linked subsurface lithology, hydraulic head variations, and vertical surface displacement to compaction rates across the six 50~m depth intervals (\Cref{fig:preprocessing_workflow}).

% ===================================================================
% SECTION 2 FLOATS (FIGURES AND TABLES AT SECTION END)
% ===================================================================

\begin{figure}[htbp]
	\centering
	\begin{tikzpicture}[
		node distance=0.7cm and 0.35cm,
		box/.style={draw, rectangle, rounded corners=3pt, align=center, fill=blue!5, line width=0.7pt, font=\footnotesize, inner sep=4pt, minimum height=0.9cm},
		proc/.style={draw, rectangle, rounded corners=3pt, align=center, fill=orange!5, line width=0.7pt, font=\footnotesize, inner sep=4pt, minimum height=0.9cm},
		resbox/.style={draw, rectangle, rounded corners=3pt, align=center, fill=green!5, line width=0.8pt, font=\small, inner sep=5pt},
		arrow/.style={-{Stealth[scale=0.9]}, line width=0.7pt, draw=gray!80}
	]
		% Step 1: Raw Inputs
		\node (raw_mlcw) [box] {\textbf{MLCW Extensometers}\\Extensometer Rings};
		\node (raw_gwl)  [box, right=of raw_mlcw] {\textbf{Groundwater Network}\\Piezometric Wells};
		\node (raw_insar)[box, right=of raw_gwl] {\textbf{cGNSS \& InSAR}\\Satellite \& Geodetic Data};
		\node (raw_coda) [box, right=of raw_insar] {\textbf{Stratigraphy}\\3D Model \& Cores};

		% Step 2: Processing Nodes
		\node (proc_mlcw) [proc, below=of raw_mlcw] {\textbf{Sectional Differencing}\\Compaction Increments};
		\node (proc_gwl)  [proc, below=of raw_gwl] {\textbf{Monthly Differencing}\\Head Changes ($\text{dGWL}$)};
		\node (proc_insar)[proc, below=of raw_insar] {\textbf{Vector Decomposition}\\Vertical Surface Motion};
		\node (proc_coda) [proc, below=of raw_coda] {\textbf{ILR / SBP Transform}\\Log-Ratio Facies Balances};

		% Step 3: Integrated Dataset
		\node (out_matrix) [resbox, below=0.8cm of $(proc_gwl)!0.5!(proc_insar)$, minimum width=12.5cm] {\textbf{Spatiotemporally Aligned Monthly Compaction Dataset}};

		% Connectors
		\draw [arrow] (raw_mlcw) -- (proc_mlcw);
		\draw [arrow] (raw_gwl) -- (proc_gwl);
		\draw [arrow] (raw_insar) -- (proc_insar);
		\draw [arrow] (raw_coda) -- (proc_coda);

		\draw [arrow] (proc_mlcw) |- (out_matrix);
		\draw [arrow] (proc_gwl) -- (out_matrix);
		\draw [arrow] (proc_insar) -- (out_matrix);
		\draw [arrow] (proc_coda) |- (out_matrix);
	\end{tikzpicture}
	\caption{Workflow diagram of multi-source observation preprocessing and spatiotemporal alignment.}
	\label{fig:preprocessing_workflow}
\end{figure}

\begin{table}[htbp]
	\centering
	\caption{Multilayer Compaction Monitoring Well (MLCW) Locations and Attributes}
	\label{tab:mlcw_info}
	\begin{tabular}{llrrrr}
		\toprule
		\textbf{ID} & \textbf{Station} & \textbf{Lon. (\si{\degree})} & \textbf{Lat. (\si{\degree})} & \textbf{Elev. (m)} & \textbf{Bottom Depth (m)} \\
		\midrule
		YL01 & Guangfu  & 120.4025 & 23.7414 & 22 & 300.0 \\
		YL02 & Huwei    & 120.4316 & 23.7153 & 25 & 300.0 \\
		YL03 & Tuku     & 120.3898 & 23.6881 & 23 & 300.0 \\
		YL04 & Honglun  & 120.3478 & 23.6866 & 17 & 340.0 \\
		YL05 & Xiutan   & 120.3496 & 23.6589 & 14 & 300.0 \\
		\bottomrule
	\end{tabular}
\end{table}

\begin{table}[htbp]
	\centering
	\footnotesize
	\setlength{\tabcolsep}{2.5pt}
	\caption{Groundwater Level (GWL) Monitoring Stations in Yunlin County}
	\label{tab:gwl_info}
	\begin{tabular}{llcc|llcc}
		\toprule
		\textbf{Station} & \textbf{Abbrev.} & \textbf{Lon. (\si{\degree})} & \textbf{Lat. (\si{\degree})} & \textbf{Station} & \textbf{Abbrev.} & \textbf{Lon. (\si{\degree})} & \textbf{Lat. (\si{\degree})} \\
		\midrule
		Annan & ANN & 120.2488 & 23.7039 & Jiulong & JLG & 120.4310 & 23.7511 \\
		Beigang & BGG & 120.3019 & 23.5789 & Kanjiao & KJO & 120.5379 & 23.6124 \\
		Bozi & BZI & 120.1515 & 23.6335 & Kecuo & KCO & 120.3342 & 23.6266 \\
		Caicuo & CCO & 120.2193 & 23.6122 & Liuhe & LHE & 120.5627 & 23.7689 \\
		Chenguang & CGU & 120.3271 & 23.5858 & Liuzhuang & JZG & 120.4006 & 23.6344 \\
		Chukou & CKU & 120.6430 & 23.7701 & Lunzi & LZI & 120.3548 & 23.6077 \\
		Citong & CTG & 120.4969 & 23.7567 & Mingde & MDE & 120.1993 & 23.6529 \\
		Dagou & DGU & 120.2107 & 23.5662 & Pingding & PDG & 120.6407 & 23.7558 \\
		Dapi & DPI & 120.4232 & 23.6629 & Qiongbu & QBU & 120.2073 & 23.5184 \\
		Dongguang & DGG & 120.2720 & 23.6519 & Sanhe & SHE & 120.4879 & 23.6051 \\
		Donghe & DHE & 120.5694 & 23.6858 & Shangyi & SYI & 120.4519 & 23.6342 \\
		Erlun & ELU & 120.4142 & 23.7717 & Shiliu & SLU & 120.5859 & 23.7207 \\
		Fangcao & FCO & 120.3741 & 23.7184 & Shuilin & SLN & 120.2461 & 23.5731 \\
		Fengrong & FRG & 120.3110 & 23.7907 & Tianyang & TYG & 120.3090 & 23.7254 \\
		Ganghou & GHU & 120.3921 & 23.7965 & Tuku & TKU & 120.3899 & 23.6879 \\
		Guangfu & GFU & 120.4024 & 23.7414 & Wencuo & WCO & 120.5121 & 23.6577 \\
		Gukeng & GKN & 120.5668 & 23.6446 & Wutu & WTU & 120.5980 & 23.7672 \\
		Haifeng & HFG & 120.2259 & 23.7649 & Xiluo & XLO & 120.4673 & 23.7959 \\
		Haiyuan & HYN & 120.1791 & 23.7208 & Xinghua & XIH & 120.2890 & 23.7603 \\
		Hefeng & HFE & 120.2234 & 23.7390 & Xiutan & XTA & 120.3496 & 23.6588 \\
		Honglun & HLN & 120.3481 & 23.6866 & Yiwu & YWU & 120.1883 & 23.5413 \\
		Houan & HAN & 120.2349 & 23.7892 & Yongding & YDN & 120.4196 & 23.7977 \\
		Huwei & HWI & 120.4323 & 23.7142 & Yuanzhang & YZG & 120.3101 & 23.6529 \\
		Huxi & HXI & 120.5112 & 23.7221 & Zhengmin & ZMN & 120.4073 & 23.7093 \\
		Huzinei & HZN & 120.5602 & 23.7834 & Zhongshan & ZHS & 120.2443 & 23.7971 \\
		Jiaxing & JXG & 120.4595 & 23.6482 & Zhongxiao & ZXO & 120.3150 & 23.6221 \\
		Jinhu & JHU & 120.1533 & 23.5734 & Zhongxing & ZXG & 120.5891 & 23.7808 \\
		\bottomrule
	\end{tabular}
\end{table}

\begin{table}[htbp]
	\centering
	\caption{Continuous GNSS Station Coordinates and Monitoring Periods}
	\label{tab:gnss_info}
	\begin{tabular}{llrrrr}
		\toprule
		\textbf{Station} & \textbf{Code} & \textbf{Lon. (\si{\degree})} & \textbf{Lat. (\si{\degree})} & \textbf{Elev. (m)} & \textbf{Period} \\
		\midrule
		Guangfu  & GFES & 120.4025 & 23.7414 & 45.07 & 2010--2024 \\
		Huwei    & NTUH & 120.4276 & 23.7330 & 48.27 & 2014--2024 \\
		Tuku     & TKJS & 120.3898 & 23.6880 & 41.13 & 2010--2024 \\
		Honglun  & HLES & 120.3478 & 23.6866 & 37.07 & 2011--2024 \\
		Xiutan   & STES & 120.3498 & 23.6587 & 38.69 & 2018--2024 \\
		\bottomrule
	\end{tabular}
\end{table}

\begin{table}[htbp]
	\centering
	\caption{Summary of the Sentinel-1A datasets used in this study.}
	\label{tab:sentinel1_info}
	\begin{tabular}{lcc}
		\toprule
		\textbf{Specifications} & \textbf{Ascending} & \textbf{Descending} \\
		\midrule
		Relative Orbit (Path) & 69 & 105 \\
		\multicolumn{1}{l}{Acquisition Period} & \multicolumn{2}{c}{April 2016--November 2021} \\
		Number of Images      & 266 & 264 \\
		\multicolumn{1}{l}{Acquisition Mode} & \multicolumn{2}{c}{Interferometric Wide (IW)} \\
		\multicolumn{1}{l}{Polarization} & \multicolumn{2}{c}{VV} \\
		Incidence Angles & $32^{\circ}$--$38^{\circ}$ & $38^{\circ}$--$43^{\circ}$ \\
		Satellite Headings & $347.63^{\circ}$ & $192.37^{\circ}$ \\
		\bottomrule
	\end{tabular}
\end{table}

\clearpage
